\chapter{Test Codes and Ordering --- Talking to Your Doctor}
\label{app:test_codes}

\textit{Scope note: This appendix provides supplementary diagnostic-testing context. Codes, coverage, and ordering systems vary by country and payer; they are not a substitute for local laboratory, hospital, or regulatory guidance.}

\clearpage
\section{Introduction}

This appendix is designed for patients who want to have informed conversations with their doctors about which tests to order. It includes the standard LOINC codes (the universal laboratory test identifiers), CPT codes (the billing codes), and a practical guide to the blood tube colors and what they mean.

You do not need to memorize these codes. But if your doctor says ``we will order a CBC,'' knowing that ``CBC'' is CPT code 85025 and means a complete blood count with differential gives you a better understanding of what is actually being tested and what you are paying for.

\clearpage
\section{Common Blood Tests --- LOINC Codes}

LOINC (Logical Observation Identifiers Names and Codes) is the universal standard for identifying lab tests. Every lab in the world uses LOINC codes to ensure that ``glucose'' in one lab means the same thing as ``glucose'' in another lab.

\begin{longtable}{@{}lll@{}}
\toprule
\textbf{Test Name} & \textbf{LOINC Code} & \textbf{Specimen} \\
\midrule
\endhead
Albumin & 1751-7 & Serum \\
Alkaline phosphatase & 6768-6 & Serum \\
ALT & 1742-6 & Serum \\
aPTT & 14959-1 & Plasma \\
AST & 1920-8 & Serum \\
\textbf{Blood Gas} & & \\
BNP & 30934-4 & Plasma \\
BUN & 3094-0 & Serum \\
Calcium & 17861-6 & Serum \\
Calcium, ionized & 17144-7 & Blood \\
\textbf{Cardiac} & & \\
CBC with differential & 58410-2 & Blood \\
Chloride & 2075-0 & Serum \\
CO$_2$ (bicarbonate) & 2028-9 & Serum \\
\textbf{Coagulation} & & \\
\textbf{Complete Blood Count (CBC)} & & \\
Cortisol & 1751-7 & Serum \\
Creatinine & 2160-0 & Serum \\
CRP (C-reactive protein) & 30522-7 & Serum \\
D-dimer & 7738-4 & Plasma \\
Direct bilirubin & 1968-7 & Serum \\
eGFR & 48642-3 & Serum \\
\textbf{Endocrine} & & \\
ESR (sed rate) & 4537-7 & Blood \\
Fibrinogen & 3255-7 & Plasma \\
Free T3 & 3027-0 & Serum \\
Free T4 & 3026-2 & Serum \\
GGT & 2324-4 & Serum \\
Glucose & 2345-7 & Serum/Plasma \\
HbA1c & 4548-4 & Blood \\
Hematocrit & 4544-3 & Blood \\
Hemoglobin & 718-7 & Blood \\
\textbf{Inflammation} & & \\
INR & 38885-1 & Plasma \\
Lactate & 2518-9 & Blood \\
LDH & 14828-2 & Serum \\
\textbf{Liver} & & \\
Magnesium & 19221-5 & Serum \\
MCH & 786-4 & Blood \\
MCHC & 787-2 & Blood \\
MCV & 785-6 & Blood \\
\textbf{Metabolic Panel} & & \\
NT-proBNP & 33762-4 & Plasma \\
pCO$_2$ & 2021-4 & Blood \\
pH (blood gas) & 2744-1 & Blood \\
Phosphorus & 2777-1 & Serum \\
Platelet count & 777-3 & Blood \\
pO$_2$ & 2703-7 & Blood \\
Potassium & 2823-3 & Serum \\
PSA & 2857-1 & Serum \\
PT (prothrombin time) & 5902-2 & Plasma \\
RBC count & 789-8 & Blood \\
RDW & 788-0 & Blood \\
Reticulocyte count & 10405-5 & Blood \\
Sodium & 2951-2 & Serum \\
Total bilirubin & 1975-2 & Serum \\
Total protein & 2885-2 & Serum \\
Troponin I & 10839-9 & Serum \\
Troponin T & 49563-0 & Serum \\
TSH & 3016-3 & Serum \\
\textbf{Tumor Markers} & & \\
\textbf{Uric Acid} & & \\
Uric acid & 14957-5 & Serum \\
\textbf{Urinalysis} & & \\
Urinalysis (complete) & 24331-1 & Urine \\
Urine creatinine & 2161-8 & Urine \\
Urine glucose & 2346-5 & Urine \\
Urine protein & 2888-6 & Urine \\
WBC count & 6690-2 & Blood \\
\bottomrule
\end{longtable}

\clearpage
\section{Common Blood Tests --- CPT Codes}

CPT (Current Procedural Terminology) codes are what your insurance company uses to bill for your lab tests. Knowing these codes helps you understand what you are being charged for and helps you compare costs if you are paying out of pocket or shopping for the best price.

\begin{longtable}{@{}lll@{}}
\toprule
\textbf{Test Name} & \textbf{CPT Code} & \textbf{Category} \\
\midrule
\endhead
Albumin & 82040 & Chemistry \\
Alkaline phosphatase & 84075 & Chemistry \\
ALT & 84460 & Chemistry \\
ANA & 86038 & Immunology \\
Anti-CCP & 86200 & Immunology \\
Anti-dsDNA & 86200 & Immunology \\
aPTT & 85730 & Coagulation \\
AST & 84450 & Chemistry \\
Basic Metabolic Panel (BMP) & 80048 & Chemistry \\
Blood culture & 87040 & Microbiology \\
BNP & 83880 & Chemistry \\
BUN & 82948 & Chemistry \\
Calcium & 82310 & Chemistry \\
\textbf{Cardiac} & & \\
CBC with differential & 85025 & Hematology \\
CBC without differential & 85027 & Hematology \\
\textbf{Chemistry} & & \\
Chloride & 82947 & Chemistry \\
CO$_2$ & 82947 & Chemistry \\
\textbf{Coagulation} & & \\
Comprehensive Metabolic Panel (CMP) & 80053 & Chemistry \\
Cortisol & 82530 & Chemistry \\
Creatinine & 82565 & Chemistry \\
CRP (high-sensitivity) & 86141 & Immunology \\
D-dimer (quantitative) & 85379 & Coagulation \\
Direct bilirubin & 82248 & Chemistry \\
eGFR & 82947 & Chemistry \\
\textbf{Endocrine} & & \\
Fibrinogen & 85384 & Coagulation \\
Free T3 & 84481 & Endocrinology \\
Free T4 & 84439 & Endocrinology \\
GGT & 82977 & Chemistry \\
Glucose & 82947 & Chemistry \\
Gram stain & 87205 & Microbiology \\
HbA1c & 83036 & Chemistry \\
HCG (serum, quantitative) & 84702 & Endocrinology \\
\textbf{Hematology} & & \\
Hepatic Function Panel & 80076 & Chemistry \\
Hepatitis B surface Ag & 87340 & Microbiology \\
Hepatitis C Ab & 86704 & Microbiology \\
HIV Ag/Ab combo & 87389 & Microbiology \\
\textbf{Immunology} & & \\
INR & 85610 & Coagulation \\
LDH & 84480 & Chemistry \\
\textbf{Liver} & & \\
Magnesium & 83735 & Chemistry \\
NT-proBNP & 84484 & Chemistry \\
Phosphorus & 84100 & Chemistry \\
Platelet count & 85049 & Hematology \\
Potassium & 82947 & Chemistry \\
Pregnancy test (urine) & 81025 & Urinalysis \\
PSA (free) & 86316 & Tumor marker \\
PSA (total) & 86300 & Tumor marker \\
PT & 85610 & Coagulation \\
Reticulocyte count & 85044 & Hematology \\
RF (rheumatoid factor) & 86431 & Immunology \\
\textbf{Infectious Disease} & & \\
RPR & 86592 & Microbiology \\
Sodium & 82947 & Chemistry \\
Total bilirubin & 82247 & Chemistry \\
Total protein & 84160 & Chemistry \\
Troponin I & 84484 & Chemistry \\
TSH & 84443 & Endocrinology \\
\textbf{Tumor Markers} & & \\
\textbf{Urinalysis} & & \\
Urinalysis (dipstick) & 81001 & Urinalysis \\
Urinalysis (microscopic) & 81015 & Urinalysis \\
Urine culture & 87088 & Microbiology \\
WBC count & 85048 & Hematology \\
\bottomrule
\end{longtable}

\clearpage
\section{Blood Tube Colors --- What Is in Each Tube}

Every color-coded tube has a specific additive inside that prepares the blood for a particular type of test. The wrong tube gives the wrong results. This is what the phlebotomist is thinking about when they draw your blood.

\begin{longtable}{@{}lll@{}}
\toprule
\textbf{Tube Color} & \textbf{Additive Inside} & \textbf{Used For} \\
\midrule
Black (black top) & Sodium citrate (different ratio) & ESR (erythrocyte sedimentation rate) \\
Blue (light blue top) & Sodium citrate & Coagulation tests (PT/INR, aPTT, D-dimer) \\
Gold/SST (red-gray top) & Gel separator + clot activator & Chemistry, endocrinology, antibodies \\
Gray (gray top) & Sodium fluoride + potassium oxalate & Glucose, lactate (prevents glycolysis) \\
Green (dark green top) & Heparin (lithium or sodium) & Chemistry, blood gas \\
Lavender/Purple (purple top) & EDTA & CBC, HbA1c, blood typing, DNA tests \\
Light green (mint top) & Lithium heparin + gel separator & General chemistry \\
Pink (pink top) & EDTA & Blood bank (type and crossmatch) \\
Red (no additive) & None & Chemistry, serology (blood clot, serum separates out) \\
Yellow (yellow top) & ACD (acid citrate dextrose) & Molecular testing, DNA studies \\
\bottomrule
\end{longtable}

\subsubsection{What This Means For You}

\textbf{Why tube color matters:} If the wrong tube is used, the test may be inaccurate or uninterpretable. For example:
\begin{itemize}
    \item If a coagulation test (PT/INR) is drawn in the wrong tube, the result is unreliable and your warfarin dose could be adjusted incorrectly.
    \item If a glucose is drawn in a red-top tube instead of a gray-top tube, the glucose will be falsely low because the blood cells will consume the sugar during clotting.
    \item If a CBC is drawn in a green-top tube instead of a lavender-top tube, the anticoagulant will not preserve the cells properly.
\end{itemize}

You do not need to tell the phlebotomist which tube to use --- this is their job. But if you see them reach for the wrong tube, it is perfectly appropriate to ask: ``Is that the right tube for my test?''

\clearpage
\section{Panel Codes --- What You Are Actually Ordering}

Most lab tests are ordered as panels (groups of related tests) rather than individually. Here is what you are getting with each common panel.

\begin{longtable}{@{}lll@{}}
\toprule
\textbf{Panel Name} & \textbf{CPT Code} & \textbf{Tests Included} \\
\midrule
Basic Metabolic Panel (BMP) & 80048 & Glucose, BUN, creatinine, sodium, potassium, chloride, CO$_2$ \\
CBC with Differential & 85025 & WBC count, RBC, hemoglobin, hematocrit, platelets, indices, differential \\
Coagulation Panel & 85610 + 85730 & PT/INR + aPTT (often ordered together) \\
Comprehensive Metabolic Panel (CMP) & 80053 & BMP + calcium, total protein, albumin, bilirubin, AST, ALT, ALP \\
Hepatic Function Panel & 80076 & Bilirubin, albumin, AST, ALT, ALP, GGT \\
Iron Studies & 83540 + 84590 + 84450 + 82040 & Iron, TIBC, transferrin saturation \\
Lipid Panel & 80061 & Total cholesterol, LDL, HDL, triglycerides \\
Thyroid Panel & 80048 + add-ons & TSH + Free T4 (sometimes Free T3) \\
Urinalysis (complete) & 81001 + 81015 & Dipstick + microscopic examination \\
\bottomrule
\end{longtable}

\clearpage
\section{How to Talk to Your Doctor About Lab Tests}

You do not need a medical degree to have a productive conversation about lab tests. Here is a framework.

\subsection{What Preparation Is Needed}

\begin{enumerate}
    \item \textbf{Write down what symptoms you have and when they started.} This helps your doctor decide which tests to order.
    \item \textbf{Make a list of all your medications, supplements, and vitamins.} This affects which tests are accurate and which may be interfered with.
    \item \textbf{Know your family medical history.} If your parents had diabetes, heart disease, or thyroid problems, you may need earlier screening.
    \item \textbf{Ask yourself what you want to know.} ``Do I want to check for diabetes?'' ``Do I want to understand why I am tired?'' ``Do I want to screen for heart disease risk?''
\end{enumerate}

\subsection{Understanding Your Results}

\begin{itemize}
    \item ``What test are you ordering and what will it tell us?''
    \item ``Do I need to fast for this test?''
    \item ``Are any of my medications likely to affect the results?''
    \item ``Should I stop any supplements before the test?''
    \item ``When will the results be available?''
    \item ``What happens if the results are abnormal?''
    \item ``How much will this test cost with my insurance?''
    \item ``Is there a less expensive alternative test that gives the same information?''
\end{itemize}

\subsection{Understanding Your Results}

\begin{itemize}
    \item ``Is this result normal for someone my age and sex?''
    \item ``How does this compare to my previous results?''
    \item ``Is this result trending in a concerning direction?''
    \item ``Do I need any follow-up testing?''
    \item ``Do I need treatment, or should we recheck this in a few months?''
    \item ``Could anything have affected this result (food, exercise, medication)?''
\end{itemize}

\clearpage
\section{Understanding Your Lab Bill}

Lab billing can be confusing. Here is what to know:

\begin{itemize}
    \item \textbf{Each test has a CPT code.} If you see multiple CPT codes on your bill, you are being charged for each test separately.
    \item \textbf{Panels are bundled.} A ``Comprehensive Metabolic Panel'' (CPT 80053) is one charge that includes 14 individual tests. You should not be billed separately for glucose, sodium, potassium, etc., if a CMP was ordered.
    \item \textbf{If you are paying out of pocket, ask for the cash price.} Lab cash prices are often dramatically lower than billed rates. Quest Diagnostics and LabCorp both publish self-pay pricing.
    \item \textbf{If your claim is denied, ask for an itemized bill} and verify that the correct CPT codes were used. Insurance denials are often due to coding errors.
    \item \textbf{Preventive screening tests are often covered at no cost} under the Affordable Care Act. This includes cholesterol screening, diabetes screening, and certain other tests as recommended by the USPSTF.
\end{itemize}
